\chapter{Zadatak 3 --- što je točno zadano}
\label{ch:z3-zadatak}

\section{Izvorni tekst, doslovno}
\label{sec:z3-izvorni}

Zadatak 3 se u izvorniku proteže preko dvije stranice. Slijede obje, u
izvornom rasporedu.

\begin{zadatak}{Zadatak 3 (Formulacija opt.\ problema; opt.\ aproksimacija s ograničenjima) --- prva stranica}
\skenzad[0.96]{z3-original-a}
\end{zadatak}

\begin{zadatak}{Zadatak 3 --- druga stranica: upute za diskretizaciju i podzadaci}
\skenzad[0.96]{z3-original-b}
\end{zadatak}

\subsection{Prijepis}

\begin{zadatak}{Zadatak 3 --- prijepis teksta}
\textbf{Zadatak 3. (Formulacija opt.\ problema; opt.\ aproksimacija s
ograničenjima)} Potrebno je izgraditi cestu na neravnom terenu. Trošak izrade
ceste proporcionalan je količini materijala (zemlje) kojeg treba oduzeti
(otkopati) ili dodati da bi se dobio donekle ,,poravnan`` teren prikladan za
izgradnju ceste.

Neka je $L$ sveukupna duljina ceste (točnije projekcija duljine ceste na
horizontalnu os $x$), a $h(x)$ je poznata visina terena u ovisnosti o koordinati
$x \in [0, L]$ --- vidjeti sliku za ilustraciju. Potrebno je odrediti visinu
ceste $y(x)$ za $x \in [0, L]$ koja rezultira minimalnim troškovima izrade, s
time da je nužno poštovati sljedeća ograničenja:
\begin{itemize}[nosep,leftmargin=*]
  \item nagib ceste ne smije biti veći od $b_1$;
  \item za ugodnu vožnju brzina promjene nagiba ceste ne smije biti veća od $b_2$;
  \item rubni uvjeti $y(0) = a$ i $y(L) = b$ moraju biti zadovoljeni.
\end{itemize}

Zadano je:
\begin{align}
L &= 10, \tag{2a} \\
h(x) &= 0.0053x^4 - 0.095x^3 + 0.48x^2 - 0.3x + 1, \tag{2b} \\
a &= h(0) = 1, \quad b = h(L) = 4, \tag{2c} \\
\text{slučaj A:}& \quad b_1 = 0.5,\ b_2 = 0.5 \tag{2d} \\
\text{slučaj B:}& \quad b_1 = 0.5,\ b_2 = 0.2 \tag{2e}
\end{align}

\textbf{Upute:} Koordinate $x$ na intervalu $[0, L]$ diskretizirajte s $N+1$
točaka $\{x_0, x_1, \dots, x_N\}$, odnosno ovaj interval podijelite na $N$
podintervala od kojih je svaki duljine $\Delta h = L/N$ (uočite $x_0 = 0$,
$x_N = L$). Tada visinu ceste računajte samo u točkama diskretizacije
$\{y_0 = y(x_0),\, y_1 = y(x_1),\, \dots,\, y_N = y(x_L)\}$. Nagib ceste
(derivaciju) aproksimirajte prvom podijeljenom razlikom:
$y'(x_i) \approx \nabla y_i \coloneqq \frac{1}{\Delta x}(y_i - y_{i-1})$, a
brzinu promjene nagiba (druga derivacija) s
$y''(x_i) \approx \nabla^2 y_i \coloneqq \frac{1}{\Delta x}(\nabla y_i - \nabla y_{i-1})$.

\textbf{Zadatci:}
\begin{itemize}[nosep,leftmargin=*]
  \item Formulirajte problem kao optimizacijski problem. Kojoj klasi
        optimizacijskih problema pripada ovaj problem?
  \item Riješite problem za 4 različita slučaja: \emph{i)} za slučaj A
        (parametri zadani s (2d)) te s $N = 20$; \emph{ii)} za slučaj A s
        $N = 100$; \emph{iii)} za slučaj B (parametri zadani s (2e)) te s
        $N = 20$; \emph{iv)} za slučaj B s $N = 100$.
\end{itemize}
\end{zadatak}

\section{Što u zadatku \emph{ne} piše}
\label{sec:z3-nije-receno}

Ovaj zadatak traži najviše vlastitog modeliranja od sva četiri --- zadan je
inženjerski opis, ne matematički problem.

\begin{longtable}{@{}p{0.36\linewidth} p{0.60\linewidth}@{}}
\toprule
\textbf{Što nije rečeno} & \textbf{Kako je protumačeno i zašto} \\
\midrule
\endfirsthead
\toprule
\textbf{Što nije rečeno} & \textbf{Kako je protumačeno i zašto} \\
\midrule
\endhead
Koja je \emph{formula} za trošak. & Zadatak kaže samo da je trošak
,,proporcionalan količini materijala koji treba otkopati ili dodati``. Uzeto je
$J = \int_0^L \abs{y(x) - h(x)}\,\dd x$: apsolutna vrijednost, jer se i iskop i
nasip plaćaju, i to \emph{jednakim} ponderom. Da su ponderi različiti, funkcija
cilja bila bi po dijelovima linearna s dva koeficijenta. \\[6pt]
Širina trase. & Nije zadana. Uzeta je jedinična, $w = 1$ m, pa je volumen
brojčano jednak površini u $\mathrm{m}^2$. Konstantna širina je pozitivan faktor
i ne mijenja minimizator. \\[6pt]
Jesu li iskop i nasip jednako skupi. & Nije rečeno. Uzeto je da jesu --- to je
najjednostavnija pretpostavka konzistentna s tekstom. \\[6pt]
Mjerne jedinice. & $L = 10$ bez jedinice; uzeto je metara. Tada je $b_1$
bezdimenzijski (m/m), a $b_2$ ima jedinicu $\mathrm{m}^{-1}$. \\[6pt]
$\Delta h$ ili $\Delta x$. & Upute u istoj rečenici koriste $\Delta h = L/N$ za
duljinu podintervala, a zatim $\frac{1}{\Delta x}$ u formulama za razlike. Riječ
je o istoj veličini pod dvama imenima; ovdje se dosljedno piše $\Delta x$. \\[6pt]
Kako uključiti apsolutnu vrijednost u LP. & Ne spominje se. Uklanja se
pomoćnim varijablama $z_i$, po postupku s Vježbi 4
(\S\ref{sec:z3-l1-trik}). \\[6pt]
Treba li skalirati sumu s $\Delta x$. & Nije rečeno. Faktor $\Delta x$ dodan je
da kriterij ima dimenziju površine; kao pozitivna konstanta ne mijenja
minimizator, samo mjerilo. \\[6pt]
Što s $\nabla^2 y_i$ na rubovima. & Druga podijeljena razlika definirana je tek
za $i \geq 2$, jer traži tri uzastopna čvora. Ograničenje glatkoće stoga vrijedi
za $N-1$ točaka, ne za $N+1$. \\
\bottomrule
\end{longtable}

\begin{zamka}[Zamka: ,,nagib ne smije biti veći od $b_1$`` znači po apsolutnoj vrijednosti]
Doslovno čitanje ,,$y'(x) \leq b_1$`` dopustilo bi proizvoljno strm \emph{pad}.
Fizikalno to nema smisla --- strmina je jednako problematična u oba smjera --- pa
se uzima $\abs{y'(x)} \leq b_1$. Isto vrijedi za $b_2$. Ta pretpostavka mijenja
broj ograničenja: svako daje \emph{dva} reda u matrici, ne jedan.
\end{zamka}

\section{Karta poglavlja}
\label{sec:z3-karta}

\begin{longtable}{@{}p{0.48\linewidth} p{0.48\linewidth}@{}}
\toprule
\textbf{Korak} & \textbf{Gdje je obrađen} \\
\midrule
\endhead
(a) Kontinuirana formulacija i klasa problema & Poglavlje~\ref{ch:z3-formulacija} \\
Diskretizacija konačnim razlikama & Poglavlje~\ref{ch:z3-diskretizacija} \\
Svođenje na standardni LP & Poglavlje~\ref{ch:z3-lp} \\
(b) Četiri slučaja, slike, analiza & Poglavlje~\ref{ch:z3-rjesenja} \\
Pitanja, zamke i modelski odgovori & Poglavlje~\ref{ch:z3-obrana} \\
\bottomrule
\end{longtable}
